\subsection{Concave neighborhoods and dual norms}
\label{subsec:concave-neighbourhood-dual-norms}
\index{Concave neighborhoods!dual support norms}
\index{Convex norm!supporting balls}
\index{Plurisubharmonic function!dual support norms}

A relatively compact neighborhood $\Omega$ of the zero section of a vector bundle $E$ determines a norm on $E^*$, the support function $\xi\mapsto\max_{v\in\overline\Omega\cap E_x}|\xi(v)|$ of its fibers. If the boundary of $\Omega$ is sufficiently concave in the horizontal directions, holomorphic sections touching the boundary from inside show that this support norm has plurisubharmonic square, and supporting balls give it uniform convexity in the fibers. Consequently, when the dual norms are obstructed as in \cref{thm:convex-norm-deformation-obstruction}, such neighborhoods cannot exist. In terms of exhaustion functions, this bounds the convexity of the complement of the zero section in the projective completion of $E$.

In this subsection, $E\to X$ is a holomorphic vector bundle of rank $r\geq1$, and $n\geq1$. Fix a smooth Hermitian metric $k$ on $E$, with real inner product $k_{\mathbb R}=\operatorname{Re}k$ and dual norm $|\cdot|_{k^*}$. For a smooth real function $f$, write $\mathcal L_f$ for its complex Hessian. Let $\Omega$ be an open set with smooth boundary, defined near $\partial\Omega$ by $f<0$ with $\mathrm df\ne0$ on $\partial\Omega$. We say that $\partial\Omega$ has \emph{at least $n$ negative Levi directions} if, at every boundary point, $\mathcal L_f$ has at least $n$ negative eigenvalues on $\ker\partial f$.

\Needspace{7\baselineskip}
\begin{lemma}\label{lem:ball-support-uniform-convexity}
\index{Convex norm!uniform strong convexity}
Fix $\delta,B>0$. On a finite-dimensional Hermitian vector space, put
\[
 s_c(\xi)=|\xi(c)|+\delta|\xi|
\]
for covectors $\xi$ and vectors $c$ with $|c|\leq B$. There is $\gamma>0$, depending only on $\delta$ and $B$, such that $s_c^2-\gamma|\xi|^2$ is real convex for every such $c$.
\end{lemma}
\begin{proof}
For $0<\tau\leq1$, replace $s_c$ by the norm
\[
 s_{c,\tau}(\xi)=\sqrt{|\xi(c)|^2+\tau|\xi|^2}+\delta|\xi|,
\]
which is smooth off zero. Its Lipschitz constant is at most $L=\sqrt{B^2+1}+\delta$, uniformly in $c$ and $\tau$. All derivatives below are real derivatives on the dual space. At $\xi\ne0$, set $u=\xi/|\xi|$ and decompose a vector as $w=\alpha u+b$ with $b$ orthogonal to $u$ for the real inner product. Write $s=s_{c,\tau}$ and $A=\mathrm ds_\xi(w)$. Since the first summand of $s$ is convex and $s\geq\delta|\xi|$,
\[
 D^2s(w,w)\geq\delta\frac{|b|^2}{|\xi|},\qquad
 D^2(s^2)(w,w)\geq2A^2+2\delta^2|b|^2.
\]
Homogeneity gives $\mathrm ds_\xi(u)=s(\xi)/|\xi|\geq\delta$. Since $|\mathrm ds_\xi(b)|\leq L|b|$, we have
\[
 |w|^2\leq\frac{2A^2}{\delta^2}
       +\left(1+\frac{2L^2}{\delta^2}\right)|b|^2.
\]
Take $\gamma=\delta^4/(\delta^2+2L^2)$. Then $\gamma\leq\delta^2/2$ and $\gamma(1+2L^2/\delta^2)=\delta^2$, so $D^2(s^2)(w,w)\geq2\gamma|w|^2$.

The homogeneous function $s^2-\gamma|\xi|^2$ of degree two extends as a $C^1$ function at zero with differential zero. Its restriction to every real line is convex, including at zero, because its second derivative is nonnegative away from zero and its first derivatives match there. Thus it is convex on the whole space. Finally,
\[
 0\leq s_{c,\tau}(\xi)-s_c(\xi)\leq\sqrt\tau\,|\xi|,
\]
and the locally uniform limit as $\tau\downarrow0$ proves the assertion with the same $\gamma$.
\end{proof}

\Needspace{10\baselineskip}
\begin{theorem}\label{thm:concave-neighbourhood-dual-norm}
\index{Concave neighborhoods!smooth dual norms}
Suppose that a relatively compact neighborhood $\Omega\Subset E$ of the zero section has smooth boundary with at least $n$ negative Levi directions. Then
\[
 p(x,\xi)=\max_{v\in\overline\Omega\cap E_x}|\xi(v)|
\]
is a continuous complex norm on $E^*$ whose square $Q=p^2$ is plurisubharmonic on the total space of $E^*$. For some $\gamma>0$, each function $Q(x,\cdot)-\gamma|\cdot|_{k^*}^2$ is real convex. If $E$ is ample, $E^*$ also admits a smooth complex norm with strictly plurisubharmonic square off the zero section.
\end{theorem}
\begin{proof}
\emph{Step 1: basic bounds.} Put $A_x=\overline\Omega\cap E_x$. By compactness of $X$ there are $0<a_0<b_0$ such that, in every fiber, the closed $k$-ball of radius $a_0$ is contained in $\Omega$ and $\overline\Omega$ lies in the open ball of radius $b_0$. The maximum defining $p$ is attained, and
\[
 a_0|\xi|_{k^*}\leq p(x,\xi)\leq b_0|\xi|_{k^*}.
\]
As a supremum of absolute values of linear functionals, $p$ is complex homogeneous and satisfies the triangle inequality.

\emph{Step 2: support points.} Fix $\xi_0\ne0$ and a maximizing point $v_0\in A_{x_0}$. It lies on $\partial\Omega$, since at an interior point a small vertical displacement would increase $|\xi_0(v)|$. Multiply $\xi_0$ by a complex number of modulus one to obtain $\eta$ with $\eta(v_0)=p(x_0,\xi_0)=a>0$. Then $\ell=\operatorname{Re}\eta$ satisfies $\ell\leq a$ on $A_{x_0}$. Let $f$ be a defining function as in the definition of negative Levi directions, so that $\Omega=\{f<0\}$ near $\partial\Omega$ and $\mathrm df\ne0$ on $\partial\Omega$.

Let $V=T^{1,0}_{v_0}E_{x_0}$ be the vertical tangent space and $T=\ker\partial f(v_0)$. Choose an $n$-dimensional complex subspace $W\subset T$ on which $\mathcal L_f$ is negative definite. We first show that the vertical differential $\mathrm d_vf(v_0)$ is nonzero. Otherwise $V\subset T$, and counting dimensions gives $W\cap V\ne\{0\}$. A vector in this intersection has negative vertical Levi form, and the two corresponding real directions show that the real vertical Hessian of $f$ has a negative direction. Its negative cone is open, hence contains a direction $w$ with $\ell(w)>0$. Since the vertical first derivative vanishes, Taylor expansion gives $f(v_0+tw)<0$ for small $t>0$, contradicting $\ell\leq a$ on $A_{x_0}$.

The boundary of $A_{x_0}$ is therefore smooth at $v_0$. Since the supporting functional $\ell$ attains its maximum on $A_{x_0}$ there,
\[
 \mathrm d_vf(v_0)=c\,\mathrm d\ell,\qquad c>0;
\]
the sign follows because an inward direction along which $\ell$ increases would contradict the maximality. On the affine hyperplane $\{\ell=a\}$ of the fiber, the restriction of $f$ has a local minimum zero at $v_0$: a point there with $f<0$ would allow an additional displacement increasing $\ell$ within $\Omega$. The real vertical Hessian of $f$ is thus nonnegative on $\ker\mathrm d\ell$. Consequently $\mathcal L_f$ is nonnegative on the complex space $V\cap T$, of dimension $r-1$. Hence $W\cap V=\{0\}$, and the projection $W\to T^{1,0}_{x_0}X$ is an isomorphism.

\emph{Step 3: plurisubharmonicity.} In local coordinates $z$ centered at $x_0$, choose a holomorphic section $\sigma_0$ through $v_0$ with tangent space $W$. The function $f\circ\sigma_0$ has zero first derivative and negative definite mixed Hessian. Its quadratic expansion is
\[
 f(\sigma_0(z))=\sum H_{\alpha\bar\beta}z_\alpha\bar z_\beta
                    +2\operatorname{Re}P(z)+O(|z|^3),
\]
where $P$ is a holomorphic quadratic polynomial. Choose a vertical vector $b$ with $\partial_vf(b)=1$, and replace $\sigma_0$ by $\sigma=\sigma_0-bP$. This preserves the first jet and cancels the pure quadratic part. Thus, after shrinking the chart,
\[
 f(\sigma(z))\leq-\epsilon|z|^2,
 \qquad \sigma(0)=v_0,
\]
for some $\epsilon>0$. In particular $\sigma(z)\in A_z$ throughout the chart.

Compactness of $\overline\Omega$ makes $p$ upper semicontinuous: for convergent covectors, the corresponding maximizing vectors have a subsequence whose limit lies in the limiting fiber of $\overline\Omega$. The supporting section gives the opposite inequality $p(z,\xi)\geq|\xi(\sigma(z))|$, with equality at $(x_0,\xi_0)$. Therefore $p$ is continuous off the zero section, and its upper bound gives continuity at zero. The holomorphic pairing $\xi(\sigma(z))$ is nonzero near the contact point, and its logarithmic absolute value is a pluriharmonic function lying below $\log p$ and equal to it at $(x_0,\xi_0)$. The mean-value property on small holomorphic disks therefore gives the submean inequality for $\log p$ at every point off the zero section. Thus $\log p$, and then $Q=\exp(2\log p)$ by Jensen's inequality, are plurisubharmonic there. At the zero section the continuous nonnegative function $Q$ vanishes, so its submean inequality on disks also holds there.

\emph{Step 4: uniform convexity in the fibers.} Consider all maximizing triples $(x,\xi,v)$ with $|\xi|_{k^*}=1$. They form a compact set, since $p$ is continuous, and by Step~2 the vertical differential $\mathrm d_vf$ does not vanish at these support points. Thus there are uniform constants $m>0$, $\rho>0$ and $C\geq1$ such that, for vertical vectors $w$,
\[
 |\mathrm d_vf|\geq m,\qquad
 f(v+w)\leq\mathrm d_vf(w)+C|w|_k^2\quad (|w|_k\leq\rho).
\]
Choose $0<\delta<\min\{a_0/2,\rho/2,m/(2C)\}$. At a support point, write $N$ for the outward real vertical unit normal and set $c=v-\delta N$. If $|w+\delta N|_k\leq\delta$, then
\[
 k_{\mathbb R}(w,N)\leq-\frac{|w|_k^2}{2\delta},\qquad
 f(v+w)\leq-\left(\frac{m}{2\delta}-C\right)|w|_k^2\leq0.
\]
All these points lie in the neighborhood where $f$ is defined, so the closed ball $\overline B_k(c,\delta)$ is contained in $A_x$.

The support function of this ball is $s_c(\xi)=|\xi(c)|+\delta|\xi|_{k^*}\leq p(x,\xi)$. Let $\xi$ be a covector of unit norm for which $v$ is a maximizing point, and let $\eta$ be $\xi$ with its phase normalized as in Step~2. Then $\operatorname{Re}\eta$ has real dual norm one, and its metric dual vector is $N$. Hence $\operatorname{Re}\eta(N)=1$, and since $|\eta(N)|\leq1$, also $\eta(N)=1$. We obtain $\eta(c)=a-\delta>0$, and therefore $s_c(\xi)=a$. Homogeneity now gives
\[
 Q(x,\xi)=\sup_c s_c(\xi)^2,
\]
where $c$ ranges over the centers of the supporting balls in the fiber. Every such center satisfies $|c|_k\leq b_0+\delta$, so by \cref{lem:ball-support-uniform-convexity} all the functions $s_c^2-\gamma|\xi|_{k^*}^2$ are convex for the same $\gamma>0$. Their supremum, which is $Q(x,\cdot)-\gamma|\cdot|_{k^*}^2$, is finite and therefore convex, as asserted.

\emph{Step 5: smoothing.} Finally, if $E$ is ample, \cref{prop:ample-dual-strict-homogeneous-function} gives the smooth strictly plurisubharmonic homogeneous function on $E^*$ required in \cref{prop:weak-homogeneous-convex-replacement}, and applying that proposition to $Q$ gives the smooth norm.
\end{proof}

For a complex manifold of dimension $N$ and an integer $1\leq q\leq N$, we say that the manifold is \emph{$q$-convex} if it has a smooth real exhaustion function with at least $N-q+1$ positive Levi eigenvalues outside a compact set; here an exhaustion function is one whose sublevel sets are compact. A $q$-convex manifold is also $q'$-convex for $q\leq q'\leq N$. We call such integers $q$ the \emph{smooth convexity indices} of the manifold, and the smallest of them, when one exists, its \emph{least smooth convexity index}.\index{$q$-convexity!least smooth convexity index}

\Needspace{9\baselineskip}
\begin{lemma}\label{lem:closed-submanifold-complement-index-bound}
\index{$q$-convexity!closed submanifold restriction}
Let $M$ be a complex manifold of dimension $N$, and let $S\subset M$ be a closed complex submanifold of dimension $D\geq1$. If $S$ has least smooth convexity index $\kappa$, then every smooth convexity index of $M$ is at least $\kappa$.
\end{lemma}
\begin{proof}
Suppose that $M$ is $p$-convex for an integer $1\leq p<\kappa$. Then a smooth exhaustion function $\psi$ has at least $N-p+1$ positive Levi eigenvalues outside a compact set $K\subset M$. Its restriction to the closed submanifold $S$ is an exhaustion function, since its sublevel sets are closed subsets of the compact sublevel sets of $\psi$, and $K\cap S$ is compact.

At a point of $S\setminus K$, choose a complex subspace $P\subset T^{1,0}M$ of dimension $N-p+1$ on which the Levi form of $\psi$ is positive definite. Then
\[
 \dim_{\mathbb C}(P\cap T^{1,0}S)\geq (N-p+1)+D-N=D-p+1.
\]
The Levi form of $\psi|_S$ is the restriction of the Levi form of $\psi$, since the inclusion is holomorphic, so it has at least $D-p+1$ positive eigenvalues. Since $p<\kappa\leq D$, this shows that $S$ is $p$-convex, which contradicts the definition of $\kappa$.
\end{proof}

In the following results, projectivizations parametrize complex lines.

\Needspace{9\baselineskip}
\begin{corollary}\label{cor:dual-norm-complement-convexity-obstruction}
\index{Projective bundle!complement convexity obstruction}
Suppose that $E$ is ample and that $E^*$ admits no smooth complex norm with plurisubharmonic square off the zero section. Then the complement
\[
 Z\setminus Y,\qquad Z=\mathbb P(\mathcal O_X\oplus E),
 \quad Y=\mathbb P(\mathcal O_X),
\]
is not $r$-convex.
\end{corollary}
\begin{proof}
The chart of $Z$ where the $\mathcal O_X$-coordinate is nonzero is identified with the total space of $E$, and $Y$ is its zero section. Suppose that $Z\setminus Y$ is $r$-convex. Since $\dim Z=n+r$, a smooth exhaustion function $\psi$ then has at least $n+1$ positive Levi eigenvalues outside a compact set $K$. Choose a relatively compact disk-bundle neighborhood $D\Subset E$ of $Y$. The set $Z\setminus D$ is compact and disjoint from $Y$. Choose a regular value $c$ of $\psi$ greater than its maximum on $Z\setminus D$ and on $K$.

Since $\{\psi\leq c\}$ is compact in $Z\setminus Y$, the set $\Omega=\{\psi>c\}\cup Y$ is an open neighborhood of $Y$. Its closure is a compact subset of $D$, and its boundary is the regular level set $\{\psi=c\}$. The defining function $f=c-\psi$ has at least $n+1$ negative Levi eigenvalues there. Restriction to the complex tangent hyperplane removes at most one negative direction, leaving at least $n$. Since $E$ is ample, \cref{thm:concave-neighbourhood-dual-norm} produces a smooth complex norm on $E^*$ with plurisubharmonic square, a contradiction.
\end{proof}

\Needspace{10\baselineskip}
\begin{corollary}\label{cor:hirzebruch-non-two-convex-complement}
\index{Hirzebruch surface!non-two-convex complements}
For every $e\geq0$, there is an ample rank-two bundle $E$ on the Hirzebruch surface $X=\mathbb F_e$ such that $\mathbb P(\mathcal O_X\oplus E)\setminus\mathbb P(\mathcal O_X)$ is not $2$-convex. The bundle may be chosen with the slope-stability and Chern classes of \cref{thm:hirzebruch-stable-ample-approximation}.
\end{corollary}
\begin{proof}
Use \cref{thm:hirzebruch-stable-ample-approximation} to choose $E$ whose dual Dolbeault operator lies in the $C^4$ neighborhood of \cref{thm:convex-norm-deformation-obstruction}; the split dual bundle satisfies the hypotheses of that theorem by \cref{prop:hirzebruch-positive-dual-tests}. Thus $E^*$ admits no smooth complex norm with plurisubharmonic square, and \cref{cor:dual-norm-complement-convexity-obstruction} with $r=2$ gives the assertion.
\end{proof}
